\section{Introduction}
\label{sec:intro}

Continuous chain-of-thought (CoT) models feed a latent vector back into the
residual stream instead of decoding an explicit reasoning token
\citep{hao2024coconut}. Matrix-valued variants of this idea introduce a
structural observable that a vector latent does not have: the \emph{rank}
of the latent matrix $\Z$. If $\Z$ carries multiple reasoning paths in
superposition, rank should track how many paths are held, and truncating
$\Z$ to rank $k$ should hurt accuracy once $k$ falls below what the task
needs.

Our companion paper, \emph{The Gradient Does Not See Rank}
\citep{larson2026gradient}, tested this directly on a matrix-CODI model (a
$d\times d$ bottleneck bolted onto a vector-pretrained GPT-2, trained by
CODI-style distillation on ProsQA) and found the probe uninformative: the
rank-$k$ truncation curve is flat to within $0.6$ percentage points across
four training regimes and four readouts, including three explicitly
nonlinear-in-$\Z$ readouts designed to rule out a constant-Jacobian
artifact. Three training seeds reached accuracies of $81.5\pm1.2$pp while
their final effective ranks spanned $\{4,12,13\}$ -- the loss rewarded no
particular rank. That paper closed the readout-mechanism question but left
one confound explicitly open (its \S7): ProsQA has a unique positive answer
and a single distractor, so the flat curves are also consistent with the
task being \emph{rank-1-solvable} -- in which case rank-blindness says
nothing about the gradient's ability to use rank when a task actually
requires it.

This paper closes that confound and asks the question the negative result
could not answer on its own:

\begin{quote}
\emph{When a task's exact solution provably requires $\rank(Z)\geq K$, does
gradient descent, trained from scratch under a hard matrix bottleneck, (a)
develop effective rank $\approx K$, and (b) make that rank causally
necessary?}
\end{quote}

Answering this honestly requires removing every shortcut a model could use
to appear rank-aware without actually needing rank: an argmax readout that
tolerates interference (\S\ref{sec:setup}), a decoder that can read raw
inputs instead of only the bottleneck (\S\ref{sec:setup}), and a fixed step
budget that mistakes a slow transition for a wall (\S\ref{sec:frontier}).
We built three experiments, each pre-registered with decision criteria
before any GPU ran, that close these shortcuts in turn.

\paragraph{Contributions.}
\begin{enumerate}
\setlength\itemsep{2pt}
\item \textbf{A provable $\rank(Z)\geq K$ lower bound with a readout
      discipline that keeps it valid.} Exact recovery of $K$ independent
      key$\to$value bindings from a linear memory needs rank $K$ -- a
      classical fact \citep{kohonen1972correlation,anderson1972simple} --
      but it holds only under \emph{exact continuous} recovery. Under
      discrete argmax decoding, a rank-$m$ matrix recovers $\approx md$
      associations \citep{nichani2025factual}, and the bound silently
      collapses. \S\ref{sec:setup} pins the readout to the literal unbind
      $\mathrm{pred}=Z\!\cdot\!\mathrm{key}_j$ scored by an absolute cosine
      bar, and closes the position-decomposition escape (any full-attention
      model can hold $K$ items at $K$ rank-1 positions) with a single-state
      $P{=}1$ bottleneck verified by a gradient-based blank-out test.
\item \textbf{Rank tracks $K$ and is causally necessary}
      (\S\ref{sec:taskd}). At $d=16$, learned effective rank tracks $K$
      almost exactly (Spearman $\rho=1.0$), and train-time force-rank
      training shows a step at $k\approx K$: at $d=8,K=4$, force-rank
      $\leq 3$ gives $0.0$ recovery and force-rank $4$ gives $0.97$.
\item \textbf{Exact composition to depth $h=21$, and the invariant-subspace
      mechanism that explains it} (\S\ref{sec:taske}). A rank-sufficient
      trained operator composes exactly under 21-fold self-application; an
      entity-subspace decomposition shows the operator restricted to the
      $K$-dimensional key subspace is exactly a scaled $K$-cycle with
      restricted rank exactly $K$, and the apparent whole-matrix
      rank-inflation to $\approx d$ is a dynamically inert complement no
      real query ever visits.
\item \textbf{Depth as a spectral-exactness amplifier.} A rank-$(K{-}1)$
      operator looks correct at shallow hops (cosine $\approx 0.94$,
      comfortably above the $0.9$ recovery bar) and is exposed only by
      21-fold composition (recovery collapses to $0.06$) -- a sharper rank
      test than any single-hop tolerance, confirmed by predicting the
      entire depth-decay curve from the operator's own eigenspectrum, no
      raw keys needed.
\item \textbf{The exactness frontier, and a methodological finding about
      declaring cells dead} (\S\ref{sec:frontier}). Scaling $d$ at fixed
      encoder width does not hit a trainability wall -- every tested seed
      eventually transitions -- but converged exactness degrades
      monotonically with $d$. Three independent axes of this program each
      called a cell ``dead'' at a fixed budget before a
      $2$--$2.5\times$ extension reversed the verdict; we name this the
      \emph{three-budget-artifacts pattern} and give the counter-example
      (a genuine, budget-independent plateau) that keeps the rule from
      being a license to always re-test.
\end{enumerate}

\paragraph{Scope, stated up front.} Every experiment here is synthetic and
$\leq$1M parameters -- deliberately small, chosen so that the ground-truth
minimal solution is provable rather than assumed. \S\ref{sec:discussion}
argues why that scope is still informative and what it does not show. We do
not claim these results demonstrate that rank helps \emph{real} reasoning;
that transfer step is explicitly future work. A brief, clearly-marked
preliminary result on a production fast-weight architecture (DeltaNet) is
reported only to bound what is already visible in a concurrent, unfinished
campaign (\S\ref{sec:discussion}); it is not this paper's headline result.
